% ============ 第八章 §5 数域和它的局部化 (原书 pp.300-308) ============
\section{数域和它的局部化}\label{sec:ch8-s5}

我们已经知道有理数域 $\Q$ 的全部赋值: $\Q$ 的全部素除子（赋值等价类）是所有的素数和 $\infty$, $\infty$ 对应 $\Q$ 中通常的阿基米德赋值, 而每个素除子 $p$ 对应非阿的 $p$ 进赋值. $\Q$ 对 $\infty$ 的完备化是实数域 $\R=\Q_{\infty}$, 而对每个素数 $p$, $\Q$ 对 $p$ 进赋值的完备化是 $p$ 进数域 $\Q_p$. $p$ 是离散素除子, 它的赋值环
\[
  \Z_p=\{\alpha\in\Q_p\mid v_P(\alpha)\geqslant 0\}
  =\{\alpha\in\Q_p\mid|\alpha|_p\leqslant 1\}
\]
叫作 $p$ 进整数环, $\Z_p$ 中元素叫 $p$ 进整数. $\Z_p$ 是 $\Z$ 在 $\Q_p$ 中对于 $p$ 进拓扑的闭包, 并且是紧集. 剩余类域 $\Z_p/p\Z_p$ 为有限域 $\F_p$.

本节要对任意代数数域 $K$, 决定 $K$ 的所有赋值和赋值等价类（素除子）. 由于 $K$ 的素除子在 $\Q$ 的限制均是 $\Q$ 的素除子, 所以问题相当于: $\Q$ 的每个素除子到 $K$ 都有哪些扩充？

先看 $\Q$ 的阿基米德赋值 $|\cdot|$ 到代数数域 $K$ 有哪些扩充. 设 $\varphi$ 是 $|\cdot|$ 到 $K$ 的一个扩充, 则 $\varphi$ 仍为阿基米德赋值. $K$ 对 $\varphi$ 的完备化 $K_{\varphi}$ 是完备的阿基米德赋值域, 而这种域只有 $(\R,\infty)$ 和 $(\C,\infty)$. 所以问题在于有多少种方法将 $K$ 看成是 $\R$ 或 $\C$ 的子域. 这个问题在本书一开始就解决了. 设 $K=\Q(\alpha)$, $[K:\Q]=n$, 则 $\alpha$ 在 $\Q$ 上的极小多项式为 $\Q[x]$ 中 $n$ 次不可约多项式 $f(x)$. $f(x)$ 在 $\C$ 中有 $n$ 个不同的根
\[
  \alpha_1,\cdots,\alpha_{r_1},\alpha_{r_1+1},\cdots,\alpha_n,
\]
其中 $\alpha_1,\cdots,\alpha_{r_1}$ 为实根, 后 $2r_2$ 个是 $r_2$ 对复根: $\alpha_{r_1+j}=\overline{\alpha}_{r_1+r_2+j}\ (1\leqslant j\leqslant r_2)$, $r_1+2r_2=n$（$\bar{\alpha}$ 表示 $\alpha$ 的复共轭）, 则 $K$ 有 $r_1$ 个实嵌入和 $r_2$ 对复嵌入
\[
  \sigma_i\colon K\hookrightarrow\R,\quad\sigma_i(\alpha)=\alpha_i\ (1\leqslant i\leqslant r_1),
\]
\[
  \sigma_i\colon K\hookrightarrow\C,\quad\sigma_i(\alpha)=\alpha_i\ (r_1+1\leqslant i\leqslant n),
\]
\[
  \sigma_{r_1+j}=\overline{\sigma}_{r_1+r_2+j}\ (1\leqslant j\leqslant r_2).
\]
所以 $\Q$ 中赋值 $|\cdot|$ 到 $K$ 只能有以下的扩充 $\varphi_i\ (1\leqslant i\leqslant n)$,
\[
  \varphi_i\colon K\to\R_{\geqslant 0},\qquad
  \varphi_i(\gamma)=|\sigma_i(\gamma)|\ (1\leqslant i\leqslant n)\quad(\gamma\in K).
\]
再看一下这 $n$ 个阿基米德赋值是否有等价的. 由于它们均为有理数域上 $|\cdot|$ 的扩充, 所以若 $\varphi_i$ 和 $\varphi_j$ 等价（即 $\varphi_i=\varphi_j^c$, $c>0$）, 则必然相等（$c=1$）. 现在设 $\varphi_i=\varphi_j\ (1\leqslant i,j\leqslant n)$, 则对每个有理数 $a$,
\begin{align*}
  |\alpha_i-a|&=|\sigma_i(\alpha)-\sigma_i(a)|=|\sigma_i(\alpha-a)|\\
  &=\varphi_i(\alpha-a)=\varphi_j(\alpha-a)=|\alpha_j-a|,
\end{align*}
这就表明 $\alpha_i$ 和 $\alpha_j$ 与每个有理数的距离均相等. 由于 $|\cdot|$ 是复数域上的连续函数, 而 $\Q$ 在 $\R$ 中稠密, 可知 $\alpha_i$ 和 $\alpha_j$ 与每个实数的距离均相等, 即 $\alpha_i$ 和 $\alpha_j$ 关于实数轴对称, 从而 $\alpha_i=\overline{\alpha}_j$, 于是 $\sigma_i=\overline{\sigma}_j$. 这就表明 $\varphi_1,\cdots,\varphi_{r_1+r_2}$ 是 $\Q$ 中阿基米德赋值到 $K$ 的彼此不同的扩充, 而 $\varphi_{r_1+r_2+j}=\varphi_{r_1+j}\ (1\leqslant j\leqslant r_2)$. 这就证明了如下结果.

\begin{theorem}\label{thm:8.20}
\begin{sloppypar}
设 $n$ 次代数数域 $K$ 在复数域 $\C$ 中共有 $r_1$ 个实嵌入 $\sigma_1,\cdots,\sigma_{r_1}$ 和 $r_2$ 对复嵌入 $\sigma_{r_1+1},\cdots,\sigma_{r_1+2r_2}$（$r_1+2r_2=n$）, $\sigma_{r_1+j}=\overline{\sigma}_{r_1+r_2+j}\ (1\leqslant j\leqslant r_2)$, 则 $K$ 共有 $r_1+r_2$ 个阿基米德素除子 $\infty_1,\cdots,\infty_{r_1+r_2}$. 以 $\varphi_i$ 表示 $\infty_i$ 中的赋值, 并且是 $\Q$ 中 $|\cdot|$ 的扩充, 则
\[
  \varphi_i(\alpha)=|\sigma_i(\alpha)|\qquad(\alpha\in K).
\]
$K$ 对于素除子 $\infty_i\ (1\leqslant i\leqslant r_1)$ 的完备化为 $(\R,|\cdot|)$, 而 $K$ 对于素除子 $\infty_i\ (r_1+1\leqslant i\leqslant r_1+r_2)$ 的完备化为 $(\C,|\cdot|)$.
\end{sloppypar}

称 $\infty_i\ (1\leqslant i\leqslant r_1)$ 为 $K$ 的\emph{实素除子}, $\infty_i\ (r_1+1\leqslant i\leqslant r_1+r_2)$ 为 $K$ 的\emph{复素除子}.
\end{theorem}

现在决定代数数域 $K$ 的非阿素除子. 仿照有理数域中的 $p$ 进赋值, 我们可以构作域 $K$ 的非阿赋值. 设 $O_K$ 是 $K$ 的整数环, $P$ 是 $O_K$ 的非零素理想. 对每个非零元素 $\alpha\in K$, 以 $v_P(\alpha)$ 表示理想 $\alpha O_K$ 作素理想分解时 $P$ 的指数. 而令 $v_P(0)=\infty$, 则 $v_P$ 是从 $K$ 到 $\Z\cup\{\infty\}$ 的映射, 并且满足指数赋值的 $3$ 条性质: 对于 $\alpha,\beta\in K$,
\[
  (1)\ v_P(\alpha)=\infty\ \text{当且仅当 } \alpha=0,
\]
\[
  (2)\ v_P(\alpha\beta)=v_P(\alpha)+v_P(\beta),
\]
\[
  (3)\ v_P(\alpha+\beta)\geqslant\min(v_P(\alpha),v_P(\beta)).
\]
前两条性质是显然的. 而第 $3$ 条是由于 $(\alpha+\beta)O_K\subseteq\alpha O_K+\beta O_K$, 而 $\alpha O_K+\beta O_K$ 作素理想分解时 $P$ 的指数为 $\min(v_P(\alpha),v_P(\beta))$. 于是 $v_P(\alpha+\beta)\geqslant\min(v_P(\alpha),v_P(\beta))$. $v_P$ 叫作域 $K$ 的 $P$ 进指数赋值. 对每个 $\gamma_P$, $0<\gamma_P<1$, 令 $|\alpha|_P=\gamma_P^{v_P(\alpha)}$（$|0|_P=0$）, 则 $|\cdot|_P$ 是域 $K$ 的非阿赋值, 并且对 $\gamma_P$ 的不同的选取, 这些 $P$ 进赋值是彼此等价的, 它们组成一个素除子, 这个素除子也表示成 $P$. 对于 $O_K$ 的不同非零素理想 $P$ 和 $Q$, $|\cdot|_P$ 和 $|\cdot|_Q$ 不等价, 从而给出不同的素除子. 由于 $v_P(K^{\times})=\Z$, 可知这些非阿素除子都是离散的. $K$ 中满足 $v_P(\pi)=1$ 的元素 $\pi$ 都是 $P$ 进素除子的素元（例如取 $\pi\in P-P^2$）. 赋值环 $O_{(P)}$ 和它的惟一极大理想 $M_P$ 为
\[
  O_{(P)}=\{\alpha\in K\mid v_P(\alpha)\geqslant 0\}
  =\{\alpha\in K\mid|\alpha|_P\leqslant 1\},
\]
\[
  M_P=\{\alpha\in K\mid v_P(\alpha)\geqslant 1\}
  =\{\alpha\in K\mid|\alpha|_P<1\}=\pi O_{(P)}.
\]
根据定义, $K$ 的整数环 $O_K$ 是 $O_{(P)}$ 的子环, 并且 $M_P\cap O_K=P$. 于是有域的单同态
\begin{equation}\tag{1}
  O_K/P\hookrightarrow O_{(P)}/M_P,
\end{equation}
从而 $K$ 对于 $P$ 进赋值的剩余类域（即上式右边）有子域 $O_K/P$. 设 $\Z\cap P=p\Z$, 其中 $p$ 为素数, 则 $K$ 的素除子 $P$ 在有理数域 $\Q$ 上的限制为素除子 $p$. 而素除子 $P$ 对 $p$ 的剩余类域次数 $f_P$ 为
\[
  f_P=[O_{(P)}/M_P:\Z/p\Z]\geqslant[O_K/P:\Z/p\Z]=f(P/p).
\]
另一方面, $v_P$ 和 $v_p$ 分别为 $K$ 的 $P$ 进标准指数赋值和 $\Q$ 的 $p$ 进标准指数赋值. 我们有 $v_P|_{\Q}=e_Pv_p$, 其中 $e_P$ 是素除子 $P$ 对 $p$ 的分歧指数. 若以 $e(P/p)$ 表示素理想 $P$ 对 $p$ 的分歧指数, 即 $pO_K$ 的素理想分解中 $P$ 的指数为 $e(P/p)$, 则 $v_P(p)=e(P/p)$. 而 $e_Pv_p(p)=e_P$, 这就表明 $e_P=e(P/p)$.

设
\[
  pO_K=P_1^{e_1}\cdots P_g^{e_g}
\]
是 $pO_K$ 的素理想分解式, 其中 $P_1,\cdots,P_g$ 为 $O_K$ 的不同非零素理想, $e_i=e(P_i/p)\geqslant 1\ (1\leqslant i\leqslant g)$, 则 $K$ 的素除子 $P_1,\cdots,P_g$ 彼此不同, 且均为 $\Q$ 的素除子 $p$ 的扩充. 由第七章定理 \ref{thm:7.5}, 我们有
\[
  \sum_{i=1}^{g}e_{P_i}f_{P_i}\leqslant[K:\Q].
\]
但是上面已证 $e_{P_i}=e(P_i/p)$, $f_{P_i}\geqslant f(P_i/p)$. 于是又有
\[
  \sum_{i=1}^{g}e_{P_i}f_{P_i}\geqslant\sum_{i=1}^{g}e(P_i/p)f(P_i/p)=[K:\Q].
\]
这就表明 $f_{P_i}=f(P_i/p)\ (1\leqslant i\leqslant g)$, 从而 (1) 式是域的同构, 即 $K$ 对素除子 $P$ 的剩余类域同构于 $O_K/P=\F_{p^f}$, 其中 $f=f(P/p)$. 并且由第七章定理 \ref{thm:7.5} 和 $\sum_{i=1}^{g}e_{P_i}f_{P_i}=[K:\Q]$ 可知 $P_1,\cdots,P_g$ 就是 $\Q$ 的素除子 $p$ 到 $K$ 的全部扩充. 由于 $K$ 的每个非阿素除子都是 $\Q$ 的某个非阿素除子的扩充, 这就证明了下面的结果:

\begin{theorem}\label{thm:8.21}
(1) $K$ 的所有非阿素除子即是全部 $P$ 进素除子（其中 $P$ 过 $O_K$ 的全体非零素理想）.

(2) 对每个素数 $p$, 以 $P_1,\cdots,P_g$ 表示环 $O_K$ 中理想 $pO_K$ 的全体素理想因子, 则 $\Q$ 的素除子 $p$ 到 $K$ 的全部扩充素除子为 $P_1,\cdots,P_g$. 并且素除子 $P_i$ 对 $p$ 的分歧指数和剩余类域次数就是第一章中定义的素理想 $P_i$ 对 $p$ 的分歧指数 $e(P_i/p)$ 和剩余类域次数 $f(P_i/p)$.
\end{theorem}

\begin{definition}\label{def:ch8-3}
对于代数数域 $K$ 的每个素除子 $P$, 以 $K_P$ 表示 $K$ 对于素除子 $P$ 的完备化, 叫作（$K$ 对 $P$ 的）\emph{局部域}.
\end{definition}

若 $P$ 是阿基米德素除子, 则 $P$ 为实素除子时 $K_P=\R$, 而 $P$ 为复素除子时 $K_P=\C$. 对于非阿素除子的情形, 我们有以下结果:

\begin{theorem}\label{thm:8.22}
设 $K$ 是代数数域, $n=[K:\Q]$, $K=\Q(\alpha)$, $\alpha\in O_K$, $f(x)=x^n+c_1x^{n-1}+\cdots+c_n\in\Z[x]$ 是 $\alpha$ 在 $\Q$ 上的最小多项式, 则对每个素数 $p$,

(1) $f(x)$ 在 $\Q_p$ 上分解成 $g$ 个不可约多项式之乘积:
\[
  f(x)=f_1(x)\cdots f_g(x),
\]
其中 $f_i(x)\ (1\leqslant i\leqslant g)$ 是 $\Z_p[x]$ 中不同的首 $1$ 多项式. 而 $g$ 为 $p$ 在 $O_K$ 中的不同素理想因子 $P_1,\cdots,P_g$ 的个数. 令 $n_i=\deg f_i(x)$, 则 $n_1+\cdots+n_g=n$.

(2) 记 $\alpha_i$ 为 $f_i(x)$ 在 $\Q_p$ 的代数闭包 $\Omega_p$ 中的一个根, 则 $K$ 对于素除子 $P_i\ (1\leqslant i\leqslant g)$ 的局部化为（同构于）
\[
  K_{P_i}\cong\Q_p(\alpha_i)\qquad(1\leqslant i\leqslant g),
\]
并且 $\Q_p(\alpha_i)$ 对于 $\Q_p$ 的分歧指数和剩余类域次数分别为 $e(P_i/p)$ 和 $f(P_i/p)$, 而 $e(P_i/p)f(P_i/p)=n_i$.
\end{theorem}

\begin{proof}
$f(x)$ 可以分解成 $\Q_p[x]$ 中首 $1$ 不可约多项式乘积, 所以可设 $f_1(x),\cdots,f_g(x)$ 是 $\Q_p[x]$ 中首 $1$ 不可约多项式. 由于 $f(x)$ 在 $\Q[x]$ 中不可约, 所以 $f(x)$ 在 $\Q$ 的任何扩域中都没有重根. 特别地, $\Q\subset\Q_p\subset\Omega_p$, 所以 $f(x)$ 在 $\Omega_p$ 中没有重根. 因此 $f_i(x)\ (1\leqslant i\leqslant g)$ 是 $\Q_p[x]$ 中不同的首 $1$ 不可约多项式. 将 $\Q$ 和 $\Q_p$ 的 $p$ 进赋值 $v_p$ 惟一地扩充成 $\Omega_p$ 中指数赋值 $v$. 由于 $f(x)$ 是 $\Z_p[x]$ 中首 $1$ 多项式, 可知对于 $f(x)$ 在 $\Omega_p$ 中的每个根 $\theta$, $v(\theta)\geqslant 0$（因为 $\theta^n+c_1\theta^{n-1}+\cdots+c_n=0$, $v(c_i)\geqslant 0$. 如果 $v(\theta)<0$, 则等式左边诸项只有 $v(\theta^n)$ 最小, 这不可能）. 所以 $f_i(x)$ 的系数作为一些根的初等对称函数, 其指数赋值 $\geqslant 0$, 但是这些系数属于 $\Q_p$, 从而必然属于 $\Z_p$, 即 $f_i(x)$ 均是 $\Z_p[x]$ 中首 $1$ 不可约多项式.

由于 $\Q\subset\Q_p\subset\Omega_p$, $\Omega_p$ 代数封闭, $\alpha$ 在 $\Q$ 上代数, 所以 $K=\Q(\alpha)$ 可作为 $\Omega_p$ 的子域. 设 $P$ 是 $p$ 在 $O_K$ 中的素理想因子, 则局部域 $K_P$ 包含 $K=\Q(\alpha)$ 和 $\Q_p$, 从而包含 $\Q_p(\alpha)$. 由于 $f(\alpha)=0$, 所以 $\alpha$ 在 $\Q_p$ 上的极小多项式为某个 $f_i(x)$. 以 $\alpha_i$ 表示 $f_i(x)$ 在 $\Omega_p$ 中的一个根, 则 $K_P\supseteq\Q_p(\alpha_i)$. 但是 $\Q_p(\alpha_i)$ 对于素除子 $P$ 是完备的, 所以 $K_P\cong\Q_p(\alpha_i)$.

如果 $\alpha_i$ 和 $\alpha_i'$ 是同一个 $f_i(x)$ 的两个根, 则 $K=\Q(\alpha)$ 嵌到 $L=\Q_p(\alpha_i)$ 和 $L'=\Q_p(\alpha_i')$ 中分别以 $\alpha\mapsto\alpha_i$ 和 $\alpha\mapsto\alpha_i'$ 的方式. 由于 $\alpha_i$ 和 $\alpha_i'$ 在 $\Q_p$ 上有同样的极小多项式 $f_i(x)$, $[L:\Q_p]=[L':\Q_p]=\deg f_i=n_i$. 若以 $v$ 和 $v'$ 表示 $v_P$ 到 $L$ 和 $L'$ 的（惟一）扩充, 则对每个 $K$ 中元素 $h(\alpha)$（其中 $h(x)\in\Q[x]$）,
\begin{align*}
  v(h(\alpha))&=v(h(\alpha_i))=\frac{1}{n_i}v_P\left(N_{L/\Q_p}(h(\alpha_i))\right)\\
  &=\frac{1}{n_i}v_P\left(N_{L'/\Q_p}(h(\alpha_i'))\right)
  =v'(h(\alpha_i'))=v'(h(\alpha_i)).
\end{align*}
这表明 $v$ 和 $v'$ 在 $K$ 上给出同样的赋值. 所以 $\Q$ 上 $p$ 进素除子到 $K$ 的扩充至多有 $g$ 个, 即至多是由 $\Q_p(\alpha_i)\ (1\leqslant i\leqslant g)$ 诱导出的素除子.

对每个 $\Q_p(\alpha_i)$, 我们有如下的域扩张图表:
\[
  \begin{array}{ccc}
    \Q_p(\alpha_i)=K_{P_i} & & \widehat{P}_i\\
    \mid & & \mid\\
    \Q_p & & \widehat{p}\\
    \mid & & \mid\\
    \Q & & p
  \end{array}
\]
以 $P_i$ 表示 $\Q_p(\alpha_i)$ 的赋值在 $K$ 上限制素除子, 则 $\Q_p(\alpha_i)=K_{P_i}$, 而 $P_i$ 是 $pO_K$ 在 $O_K$ 中的素理想因子. 以 $\widehat{P}_i$ 表示 $\Q_p(\alpha_i)$ 的素除子, $p$ 和 $\widehat{p}$ 表示 $\Q$ 和 $\Q_p$ 的 $p$ 进素除子. 由定理 \ref{thm:8.2} 知
\[
  e(\widehat{P}_i/p)=e(\widehat{P}_i/P_i)e(P_i/p)=e(P_i/p),
\]
\[
  e(\widehat{P}_i/p)=e(\widehat{P}_i/\widehat{p})e(\widehat{p}/p)=e(\widehat{P}_i/\widehat{p}),
\]
这就表明 $e(P_i/p)=e(\widehat{P}_i/\widehat{p})$. 同样可证 $f(P_i/p)=f(\widehat{P}_i/\widehat{p})$, 于是
\[
  \sum_{i=1}^{g}e(P_i/p)f(P_i/p)
  =\sum_{i=1}^{g}e(\widehat{P}_i/\widehat{p})f(\widehat{P}_i/\widehat{p})
  =\sum_{i=1}^{g}n_i=n.
\]
再由定理 \ref{thm:8.21} 便知 $P_i\ (1\leqslant i\leqslant g)$ 是 $K$ 的不同素除子. 从而 $f(x)$ 在 $\Q_p[x]$ 中的不可约因子个数 $g$ 等于 $pO_K$ 在 $O_K$ 中素理想因子的个数.
\end{proof}

\begin{example}
决定 $K=\Q(\sqrt[3]{2})$ 的全部素除子.

\emph{解} 以 $\omega$ 表示复数域中的 $3$ 次本原单位根, 则 $\alpha=\sqrt[3]{2}$ 的共轭元素为 $\alpha,\alpha\omega$ 和 $\alpha\omega^2$. $f(x)=x^3-2$ 是 $\alpha$ 在 $\Q$ 上的极小多项式. $1,\alpha,\alpha^2\in O_K$. 计算判别式
\begin{align*}
  \operatorname{disc}(1,\alpha,\alpha^2)&=
  \begin{vmatrix}
    1 & 1 & 1\\
    \alpha & \omega\alpha & \omega^2\alpha\\
    \alpha^2 & \omega^2\alpha^2 & \omega\alpha^2
  \end{vmatrix}^2
  =\alpha^6
  \begin{vmatrix}
    1 & 1 & 1\\
    1 & \omega & \omega^2\\
    1 & \omega^2 & \omega
  \end{vmatrix}^2\\
  &=4
  \begin{vmatrix}
    3 & 1 & 1\\
    0 & \omega & \omega^2\\
    0 & \omega^2 & \omega
  \end{vmatrix}^2
  =36
  \begin{vmatrix}
    \omega & \omega^2\\
    \omega^2 & \omega
  \end{vmatrix}^2\\
  &=36(\omega^2+\omega-2)=-108,
\end{align*}
可知域 $K$ 的判别式 $d(K)$ 是 $-108=-2^2\cdot 3^3$ 的因子, 并且 $-108/d(K)$ 是整数的平方. 可知 $d(K)$ 必有因子 $3$, 即 $3$ 在 $K$ 中分歧. 而当 $p\geqslant 5$ 时, $p$ 在 $K$ 中不分歧.

(1) 对于 $p=3$, $K=\Q(\alpha)=\Q(\alpha+1)$, 而 $\alpha+1$ 的极小多项式 $f(x-1)=(x-1)^3-2=x^3-3x^2+3x-3$ 是 $\Q_3[x]$ 中的 Eisenstein 多项式. 从而 $\Q_3(\alpha)=\Q_3(\alpha+1)$ 是 $\Q_3$ 的 $3$ 次完全分歧扩张. 因此 $3O_K=P^3$, $e(P/3)=3$, $f(P/3)=1$. $v_3$ 到 $K$ 有惟一的扩充 $v_P$. 对于 $K$ 中元素 $\gamma=a+b\alpha+c\alpha^2$（$a,b,c\in\Q$）,
\[
  v_P(\gamma)=\frac{1}{3}v_3(N_{K/\Q}(\gamma)),
\]
其中
\begin{align*}
  N_{K/\Q}(\gamma)&=(a+b\alpha+c\alpha^2)(a+b\alpha\omega+c\alpha^2\omega^2)(a+b\alpha\omega^2+c\alpha^2\omega)\\
  &=a^3+b^3+c^3-3abc.
\end{align*}

(2) 对于 $p=2$, $f(x)=x^3-2$ 是 $\Q_2[x]$ 中的 Eisenstein 多项式. 所以也有 $2O_K=P^3$（完全分歧）. $v_2$ 在 $K$ 中只有一个扩充 $v_P$,
\[
  v_P(\gamma)=\frac{1}{3}v_2(a^3+b^3+c^3-3abc).
\]

以下设 $p\geqslant 5$.

(3) 若 $p\equiv 2\ (\mathrm{mod}\ 3)$（$p\geqslant 5$）, 则 $\F_p^{\times}$ 是 $p-1$ 阶循环群. 由于 $3\nmid(p-1)$, 可知 $\F_p^{\times}$ 中每个元素都是另一元素的立方, 特别地, $2$ 是模 $p$ 的三次剩余, 即存在整数 $s$, 使得 $s^3\equiv 2\ (\mathrm{mod}\ p)$. 从而对于 $f(x)=x^3-2$,
\[
  v_p(f(s))=v_p(s^3-2)\geqslant 1>0=2\cdot v_p(3s^2)=2\cdot v_p(f'(s)).
\]
由 Hensel 引理知 $f(x)$ 在 $\Z_p$ 中有根 $\theta$. 而 $f(x)$ 在 $\Omega_p$ 中另两个根为 $\theta\omega$ 和 $\theta\omega^2$, 其中 $\omega$ 是 $\Omega_p$ 中的 $3$ 次本原单位根. 由于 $\theta\in\Q_p$, 可知
\[
  \theta\omega\in\Q_p\Leftrightarrow\omega\in\Q_p\Leftrightarrow\Q_p(\omega)=\Q_p.
\]
但是 $[\Q_p(\omega):\Q_p]$ 等于满足 $p^f\equiv 1\ (\mathrm{mod}\ 3)$ 的最小正整数 $f$. 由 $p\equiv 2\ (\mathrm{mod}\ 3)$ 可知 $f=2$, 这表明 $\omega\notin\Q_p$, 所以 $f(x)=x^3-2$ 在 $\Q_p$ 中只有一个根 $\theta$, 即 $f(x)$ 在 $\Q_p[x]$ 中分解成两个不可约多项式 $(x-\theta)$ 和 $x^2+\theta x+\theta^2$ 的乘积. 由于 $p$ 在 $K$ 中不分歧, 可知
\[
  pO_K=PQ,\qquad e(P/p)=e(Q/p)=1,\qquad f(P/p)=1,\ f(Q/p)=2,
\]
\[
  K_P=\Q_p(\theta)=\Q_p,\qquad K_Q=\Q_p(\theta\omega)\ (\text{为 } \Q_p\ \text{的二次扩域}).
\]
$v_p$ 在 $K$ 中有两个扩充 $v_P$ 和 $v_Q$. 对于 $K$ 中元素 $\gamma=a+b\alpha+c\alpha^2$（$a,b,c\in\Q$）,
\[
  v_P=v_p(a+b\theta+c\theta^2),
\]
\[
  \begin{aligned}
    v_Q&=\frac{1}{2}v_p\left((a+b\theta\omega+c\theta^2\omega^2)(a+b\theta\omega^2+c\theta^2\omega)\right)\\
    &=\frac{1}{2}v_p\left(a^2+b^2\theta^2+c^2\theta^4-ab\theta-ac\theta^2-bc\theta^3\right).
  \end{aligned}
\]

(4) 设 $p\equiv 1\ (\mathrm{mod}\ 3)$. 这时 $\Q_p$ 中有 $3$ 次本原单位根 $\omega$. 如果 $2$ 是模 $p$ 的三次剩余, 由 Hensel 引理知 $f(x)=x^3-2$ 在 $\Q_p$ 中有根 $\theta$, 从而另外两个根 $\theta\omega$, $\theta\omega^2$ 也属于 $\Q_p$. 于是在 $\Q_p[x]$ 中 $x^3-2$ 分解为 $(x-\theta)(x-\theta\omega)(x-\theta\omega^2)$. 所以
\[
  pO_K=P_1P_2P_3,\qquad e(P_i/p)=f(P_i/p)=1\quad(1\leqslant i\leqslant 3),
\]
其中 $P_1,P_2,P_3$ 是素除子 $p$ 在 $K$ 中 $3$ 个不同的扩充, $K_{P_i}=\Q_p\ (1\leqslant i\leqslant 3)$. 以 $v_{P_i}\ (1\leqslant i\leqslant 3)$ 表示 $v_p$ 在 $K$ 中 $3$ 个不同的扩充, 则对于 $K$ 中元素 $\gamma=a+b\alpha+c\alpha^2$（$a,b,c\in\Q$）,
\[
  v_{P_1}(\gamma)=v_p(a+b\theta+c\theta^2),
\]
\[
  v_{P_2}(\gamma)=v_p(a+b\theta\omega+c\theta^2\omega^2),
\]
\[
  v_{P_3}(\gamma)=v_p(a+b\theta\omega^2+c\theta^2\omega).
\]
如果 $2$ 不是模 $p$ 的三次剩余, 则 $f(x)=x^3-2$ 在 $\Q_p[x]$ 中不可约. 这时 $pO_K=P$, $e(P/p)=1$, $f(P/p)=3$, $v_p$ 在 $K$ 中有惟一扩充 $v_P$, $K_P=\Q_p(\alpha)$ 是 $\Q_p$ 的 $3$ 次扩域. 对于 $K$ 中元素 $\gamma=a+b\alpha+c\alpha^2$（$a,b,c\in\Q$）,
\[
  v_P(\gamma)=\frac{1}{3}v_p(a^3+b^3+c^3-3abc).
\]

现在决定 $K=\Q(\sqrt[3]{2})$ 的阿基米德素除子. 由于 $x^3-2$ 在复数域中有一个实根 $\alpha=\sqrt[3]{2}$ 和一对复根 $\alpha\omega,\alpha\omega^2$, 所以 $K$ 有一个实的素除子 $\infty_1$ 和一个复的素除子 $\infty_2$. $K_{\infty_1}=\R$, $K_{\infty_2}=\C$. 素除子 $\infty_1$ 和 $\infty_2$ 中分别有 $\Q$ 中通常赋值 $|\cdot|$ 的扩充 $|\cdot|_1$ 和 $|\cdot|_2$, 其中对 $K$ 中元素 $\gamma=a+b\alpha+c\alpha^2$（$a,b,c\in\Q$）,
\begin{align*}
  |\gamma|_2&=\left|(a+b\alpha\omega+c\alpha^2\omega^2)(a+b\alpha\omega^2+c\alpha^2\omega)\right|^{\frac{1}{2}}\\
  &=\left|a^2+b^2\alpha^2+c^2\alpha^4-ab\alpha-ac\alpha^2-bc\alpha^3\right|^{\frac{1}{2}}\\
  &=\left|a+b\alpha\omega+c\alpha^2\omega^2\right|.
\end{align*}
\end{example}

\begin{Exercises}
\item 对每个二次数域 $K=\Q(\sqrt{d})$, 决定 $\Q$ 的阿基米德赋值 $|\cdot|$ 和所有 $p$ 进指数赋值 $v_P$ 在 $K$ 中的全部扩充.

\item 设 $K$ 是代数数域, $p$ 为素数, $P$ 是 $pO_K$ 在 $O_K$ 中的一个素理想因子. $O_P$ 为局部域 $K_P$ 的赋值环, 即 $O_P=\{\alpha\in K_P\mid v_P(\alpha)\geqslant 0\}$. 求证 $O_P$ 是 $\Z_p$ 在 $K_P$ 中的整闭包, 即对每个 $\alpha\in K_P$,
\[
  \alpha\in O_P\Leftrightarrow\alpha\ \text{在 } \Z_p\ \text{上整}
  \left(\text{即 } \alpha\ \text{是 } \Z_p[x]\ \text{中某个首 } 1\ \text{多项式的根}\right).
\]

\item 决定 $\Q$ 的指数赋值 $v_5$ 到域 $K=\Q(\sqrt{2},\sqrt{-1})$ 的所有扩充. 决定 $\Q$ 的阿基米德赋值 $|\cdot|$ 到 $K$ 的所有扩充.

\item 设 $K$ 是 $\Q_p$ 的有限次完全分歧扩张, $\pi$ 是 $K$ 的一个素元. 求证 $K$ 的赋值环为 $\Z_p[\pi]$.

\item 设 $K$ 是 $\Q_p$ 的 $n$ 次不分歧扩张, $O$ 是 $K$ 的赋值环. 求证 $O$ 中存在 $p^n-1$ 次本原单位根 $\alpha$, 并且 $O=\Z_p[\alpha]$.
\end{Exercises}
